% !PRINTED 586
\weberpage{586}
Hierin sind $F_\nu$, $\psi_\nu$ ganze Functionale des absoluten Rationalitätsbereiches,
während $f_\nu$ ganze Functionale in $R$ sind. Die Ausdrücke für $F_\nu, f_\nu, \psi_\nu$
ergeben sich aus den Formeln \S.\,158, (15). Auf diese Functionale kann man dieselben
Schlüsse anwenden, die zu dem Theorem \S.\,158, 4. geführt haben, und es ergiebt sich:

\weq{12}{
\begin{aligned}
\mathfrak{S}_R \frac{f_\nu(\tau)}{f'(\tau)} &= 0, & \nu &= 0, 1, \ldots, n-2.\\
\mathfrak{S}_R \frac{f_{n-1}(\tau)}{f'(\tau)} &= 1.
\end{aligned}
}
Wenn man nun die Functionen $F_\nu(t)$, wenn sie von höherem als $(n-1)^{\mathrm{ten}}$
Grade sind, durch $f(t)$ dividirt und den Rest der Division nimmt; dann in diesem Reste
die Potenzen von $t$ [nach \S.\,158, (15), auf $f_\nu$ angewandt], durch die Functionen
$f_\nu(t)$ ausdrückt, so ergiebt sich die Darstellung

\weq{13}{
F_\nu(\tau) = \alpha_0 f_0 + \alpha_1 f_1 + \cdots + \alpha_{n-1} f_{n-1},
}
worin die $\alpha_0, \alpha_1, \ldots, \alpha_{n-1}$ \gesperrt{ganze Functionale} in $R$ sind. Diese
Darstellung gilt für jede Wurzel $\tau$ von $f(t) = 0$, d.\,h. für $\tau = \tau_1, \tau_2,
\ldots, \tau_n$. Hieraus ergiebt sich aber nach (12):
\[
\mathfrak{S}_R \left( \frac{F_\nu(\tau)}{f'(\tau)} \right) = \alpha_{n-1}.
\]
Dividirt man hier durch $\psi'(\sigma)$ und bildet die Spur $S_R$, so erhält man aus (7)
mit Anwendung des Satzes 4., \S.\,158:
\weq{14}{
S_\Omega \left( \frac{F_\nu(\tau)}{f'(\tau)\,\psi'(\sigma)} \right) = a_\nu,
}
worin $a_\nu$ ein ganzes rationales Functional bedeutet.

Nach (9) ist nun für jede Wurzel $\tau_i$ von $F(t) = 0$ mit Ausnahme von $\tau_i = \tau$
\[
\tau^{n\,m-1} F_0(\tau_i) + \tau^{n\,m-2} F_1(\tau_i) + \cdots + F_{n\,m-1}(\tau_i) = 0,
\]
und
\[
\tau^{n\,m-1} F_0(\tau) + \tau^{n\,m-2} F_1(\tau) + \cdots + F_{n\,m-1}(\tau) = F'(\tau),
\]
so dass sich, wenn man (14) mit $\tau^{n\,m-\nu-1}$ multiplicirt und die Summe in Bezug
auf $\nu$ nimmt,
\weq{15}{
\frac{F'(\tau)}{f'(\tau)\,\psi'(\sigma)} = a_0 \tau^{n\,m-1} + a_1 \tau^{n\,m-2} + \cdots
}
ergiebt. \gesperrt{Es ist daher $F'(\tau)$ durch $f'(\tau)\,\psi'(\sigma)$ theilbar.}
